\chapter{Prefacio}

Este libro nace de una ambición sencilla: un único curso de ciencias de la
Tierra y del espacio, coherente de principio a fin, que abarque desde el
primer año de escuela hasta el final del grado universitario, y que pueda
leerse en cualquier parte del mundo.

El estilo es conciso y riguroso: un curso construido a partir de
definiciones, ejemplos, proposiciones, métodos y esquemas, con un
razonamiento cuidadoso siempre que resulte accesible al nivel
considerado, seguido de ejercicios cuidadosamente elegidos. Los
resultados que se enuncian sin una deducción completa se señalan
explícitamente como admitidos. Las soluciones completas de todos los
ejercicios se recogen al final del libro. Cada número citado sobre la
Tierra real, sus océanos, su atmósfera o el cielo procede de una fuente de
referencia, y cada curva trazada a partir de medidas reales se calcula a
partir de los datos publicados.

El curso consta de cinco libros. El primero abarca los años de primaria y
de secundaria obligatoria, el segundo los tres años de bachillerato, y los
tres últimos los tres años de universidad. Los dos libros escolares nunca
enseñan dos veces la misma noción: cada idea se enseña una sola vez, en el
año en que encaja por primera vez, y un capítulo posterior que la necesite
empieza con un breve recordatorio, \emph{Lo que ya sabes}, antes de ir más
allá. Cuando una noción pertenece a la física, a la biología o a la
química, la enseñan los libros hermanos de la misma colección, y este libro
la recuerda antes de adoptar la mirada del especialista en ciencias de la
Tierra. En los libros escolares, cada año es una parte, dividida en
capítulos; cada libro universitario es un año, dividido en capítulos.
Cuanto más temprano es el año, más joven es el lector al que se dirige:
más ilustraciones, más pasos detallados y ejercicios más suaves.

\section*{Cómo leer este libro}

Los enunciados van numerados dentro de cada capítulo y se distinguen por
colores: las definiciones en azul, los teoremas en rojo, las
proposiciones y sus parientes en naranja, y los métodos --- recetas
breves para las tareas habituales --- en verde. Los ejercicios están
graduados desde ($\star$), de aplicación directa, hasta
($\star\star\star$), los problemas más difíciles. Los recuadros
enmarcados recogen lo que no forma parte del razonamiento: un
recordatorio de lo que enseñó un capítulo anterior, cómo trabajan los
geólogos y los demás especialistas en ciencias de la Tierra sobre el
terreno o en el laboratorio, una página de historia, cifras que cambian de
un año a otro (con el año al que corresponden) y consejos de seguridad.

\section*{Seguridad}

La Tierra se aprende mirándola: el cielo, el tiempo que hace, un paisaje,
las piedras de un edificio, la tierra de un jardín. Este libro no pide
nada más a su lector. No mires nunca al Sol, ni a simple vista ni
con un instrumento: ciega. Las canteras, los acantilados, las cuevas, los
ríos crecidos y los volcanes son lugares para profesionales con formación
y equipo, y este libro describe su trabajo sin enviar a nadie allí.

\section*{Cómo contribuir}

Las fuentes de este libro están en un repositorio git público:
\begin{center}
\url{https://github.com/bvirrion/one-earth-and-space-sciences-book}
\end{center}
Toda contribución --- nuevos capítulos, correcciones, mejores
explicaciones, ejercicios adicionales --- es bienvenida; véase
\texttt{CONTRIBUTING.md} en la raíz del repositorio.
